\lysh{21162_SOLUTION}{1}
% TODO: TikZ σχήμα (μην το κατασκευάσεις)
α) Το μέσο $M$ του ευθυγράμμου τμήματος $AB$ έχει συντεταγμένες:
\[ M\left(\frac{3+(-1)}{2}, \frac{2+(-6)}{2}\right) = M(1,-2). \]
β) Ο συντελεστής διεύθυνσης της ευθείας που διέρχεται από τα σημεία $A$ και $B$ είναι:
\[ \lambda_{AB} = \frac{2-(-6)}{3-(-1)} = \frac{8}{4} = 2. \]
γ) Επειδή η ευθεία $AB$ είναι κάθετη στην ευθεία $(\varepsilon)$, τότε $\lambda_{AB} \cdot \lambda_\varepsilon = -1$.
\[ 2 \cdot \lambda_\varepsilon = -1 \Leftrightarrow \lambda_\varepsilon = -\frac{1}{2}. \]
Επιπλέον, από το β) $\lambda_{AB} = 2$ άρα:
Επομένως η μεσοκάθετος $(\varepsilon)$ του τμήματος $AB$ έχει εξίσωση
\[ y+2 = -\frac{1}{2}(x-1) \Leftrightarrow x+2y+3 = 0. \]
\textbf{Εναλλακτική λύση}
Ένα σημείο $M(x,y)$ ανήκει στη μεσοκάθετο του ευθυγράμμου τμήματος $AB$ αν και μόνο αν
\[ d(M,A) = d(M,B) \Leftrightarrow \sqrt{(3-x)^2 + (2-y)^2} = \sqrt{(-1-x)^2 + (-6-y)^2} \]
\[ (3-x)^2 + (2-y)^2 = (1+x)^2 + (6+y)^2 \]
\[ 9-6x+x^2+4-4y+y^2 = 1+2x+x^2+36+12y+y^2 \]
\[ 9-6x+4-4y = 1+2x+36+12y \Leftrightarrow 13-6x-4y = 37+2x+12y \]
\[ 13-6x-4y = 37+2x+12y \Leftrightarrow -6x-2x-4y-12y = 37-13 \Leftrightarrow -8x-16y = 24 \]
\[ -8x-16y = 24 \Leftrightarrow x+2y = -3 \Leftrightarrow x+2y+3 = 0. \]